% =============================================================================
% Chapter 5 — Discussion
% =============================================================================

\chapter{Discussion}
\label{ch:discussion}

\section{The Claim, Tested}
\label{sec:claim-tested}

Chapter~2 staked the dissertation on a specific claim: that algorithmic governance tools are structurally dual-use — transferable between civilian urban management and military spatial destruction with a fluidity that physical infrastructure cannot match — and that this fluidity is produced by the economics of software, the institutional logic of control societies, and the epistemic authority of algorithmic systems. The systematic map in Chapter~4 lets that claim be tested against 270 coded records rather than asserted from a handful of canonical texts. The verdict is more nuanced than simple confirmation. The literature supports each element of the framework — but it supports them in different literatures, and the junction between them is nearly empty.

\subsection{Arrow: confirmed by absence}

The economic mechanism — Arrow's near-zero marginal cost of reproduction as the substrate of dual-use fluidity — is the framework's essential element and the literature's loudest silence (§\ref{sec:theme1}). Eighty-six records document transfer between civilian and military domains; almost none theorise its cost structure. The physical-dual-use counterexamples make this absence instructive: in \textcite{zeitoun2018}'s Basrah, dual-use sanctions bit on physical assets — steel tankers, spare parts, the material supply chains of state infrastructure. In \textcite[p.~8]{salamanca2011}'s Gaza, weaponised electricity required switches, fuel lines, and an ongoing military siege. Software escapes all this friction. The corpus — practitioner papers designing the transfer \parencite{tsiolis2025} no less than critical accounts deploring it \parencite{amoore2009,wall2013} — simply takes that escape for granted. The framework's first element is thus confirmed in the strongest way a literature review permits: the phenomenon is documented pervasively, its explanation is absent, and the explanation the framework provides contradicts nothing in the record.

\subsection{Deleuze and DeLanda: confirmed by convergence}

The control-society and machinic-assemblage elements fare differently: they are not absent but saturated (§\ref{sec:theme3}). Deleuze travels explicitly across policing, smart-city, and occupation cases alike. More telling is the convergence the review was designed to detect: military doctrine — DARPA's Combat Zones That See profiling the ``normal'' rhythms of ``entire subject cities'' \parencite[p.~35]{graham2008}, the ``military urban'' as system-of-systems \parencite[p.~406]{danielsson2025}, the city brain as OODA loop \parencite[pp.~3--4]{weber2023} — arrives independently at precisely the ontology that DeLanda's targetable city predicts and that smart-city practitioners advertise. When smart-city advocates, critical theorists, and defence planners share an ontology, the ontology is no longer a critical imposition; it is a description. And the analogue cases \parencite{coyles2017,crane2017} dissolve the obvious objection — that the targetable city is hype — by showing the same modulation logic operating bureaucratically decades before the dashboards. Software automates it.

\subsection{Michael: confirmed and extended}

The epistemic-authority element is the corpus's most populated theme (n=117), and the literature extends Michael in two directions he did not go. First, institutional: the black box is manufactured — by trade secrecy, procurement frameworks, and statutory carve-outs (§\ref{sec:theme2}) — which relocates epistemic authority from the general's briefing room to the vendor contract and the regulation's exemption clause. Second, constitutive: the strongest records reject the `bias' frame in favour of race as an originary property of what these systems are built to do \parencite{sharma2024,ziadah2021}. Michael's claim was that the military became the arbiter of truth when civilian politics abdicated. The corpus suggests the algorithmic version is more radical: the system does not merely arbitrate truth, it produces the categories — `person of interest,' `pattern of life,' `associated' — within which truth-claims are thinkable at all \parencite{amoore2009}.

\subsection{King, Graham, Weizman: the canon holds, the link is missing}

The destruction canon emerges intact and intensely alive — 34 of 76 Theme-4 records score 3 — and is currently undergoing its fastest expansion, driven by Gaza 2023–25 \parencite{benguita2025,holail2024}. Its key internal development is the cyclical reading: renewal produces displaceability, razing follows, reconstruction continues the war by other means \parencite{genc2021,smith2022,hourani2024}, with ``Dracula urbanism'' \parencite[p.~2]{wilson2026} extending domicide into peacetime development logic. What the canon does not yet contain is the algorithmic mode of destruction. That absence is the subject of the next section.

\section{The Empty Intersection}
\label{sec:empty-intersection}

The review's central empirical finding is negative: the intersection of the algorithmic-governance literature and the spatial-destruction literature — records that substantively treat algorithmically mediated urbicide — comprises five records (§\ref{sec:thin-edge}), all but one published since 2021. The two literatures were found by different instruments, cite different canons, and score inversely on the framework's first and last themes (Table~\ref{tab:themes}). Even the snowball's anchor structure tells the story: Kitchin's network produced the governance side, Graham's and Weizman's the destruction side, and Michael's — the theorist whose concept bridges them — produced nothing the other networks had not already captured.

Three explanations are plausible, and the corpus cannot fully adjudicate between them. The first is disciplinary path dependence: urbicide theory descends from critical geography and architecture; algorithmic-governance critique from surveillance studies and STS; their conferences, journals, and citation graphs barely touch. The second is secrecy: the systems that would constitute the intersection — Lavender, Where's Daddy?, the Gospel — are classified and contested; scholars can access them only through investigative journalism \parencite{hourani2024}, which throttles the normal machinery of theoretical elaboration. The third is timing: the large-scale deployment of AI target-generation in an urban war is, at the time of writing, roughly two years old. The literature is not absent; it is embryonic. On this reading, the empty intersection is not evidence against the framework but evidence of the framework's timeliness — and it fixes the dissertation's contribution as cartographic: naming a junction before the traffic arrives.

\section{Dual-Use as Warrant}
\label{sec:dual-use-warrant}

One finding exceeded the framework, and it deserves emphasis because it inverts the usual direction of critique. The framework treats dual-use as an analytic category — a lens for seeing how algorithmic systems travel between domains. But the corpus shows dual-use operating as a juridical category inside the machinery of destruction: hospitals, schools, mosques, and universities ``struck under the rationale of dual-use'' \parencite[p.~406]{benguita2025}; ``dual use facilities'' as the sanction-era targeting warrant in Iraq \parencite[p.~406]{coward2009}; dual-use electrical infrastructure as the targeting category whose deliberate degradation cost up to 100,000 civilian lives \parencite[p.~355]{graham2011}. The Iraqi power grid was not destroyed incidentally; it was targeted specifically because its dual-use character — electricity serves both civilian hospitals and air-defence radars — made it militarily licit under prevailing interpretations of IHL. Gaza 2023–25 operates the same warrant on a vastly expanded scale, where the civilian data infrastructure itself — population maps, building footprints, connectivity graphs — becomes both the targeting instrument and the targeting justification. The category that export-control law uses to restrain the flow of sensitive technology is the same category targeting doctrine uses to license the destruction of civilian fabric. This is not a contradiction; it is the same operation seen from two ends. Dual-use names the collapse of the civilian/military distinction — and that collapse is as much the targeter's resource as the critic's concern.

The warrant logic is not confined to targeting doctrine; it is stated, by at least one state, as national strategy. Israel's \emph{National Initiative for Secured Intelligent Systems} — the 2020 Special Report to the Prime Minister authored by \textcite{benisrael2020} — explicitly frames intelligent systems as a strategic technology whose national value lies in the dual combination of security needs and economic opportunity, and places that duality at the centre of state planning \parencite{benisrael2020}. This is the dual-use logic in the language of national doctrine, from the state whose military would, within three years, be deploying AI-driven target-generation systems in an urban war. The convergence of strategic policy discourse and operational military practice is not coincidental: it is the structural dual-use of this dissertation's framework, formalised at the highest level of state planning. Any policy response that strengthens the dual-use regime without reckoning with its targeting function will, at best, police the supply side of a logic it leaves intact at the point of use.

\section{Policy Implications}
\label{sec:policy-implications}

For a public-policy dissertation, the map carries four implications, escalating in specificity.

First, \textit{regulate the ontology, not the artefact}. Current dual-use regimes enumerate items — sensors, chips, software categories. The corpus shows transfer moving through data models, ontologies, and interoperability standards (§\ref{sec:theme1}), and through vendor contracts that carry military logics into civilian administration wholesale (Palantir's counterinsurgency-to-LAPD pipeline; the HART tool built on Experian marketing segments). A governance regime that licences artefacts while leaving the ontological layer unexamined will be perpetually one step behind the transfer it purports to control.

Second, \textit{opacity is a policy output and can be a policy target}. If the black box is manufactured by procurement and statute (§\ref{sec:claim-tested}), it can be unmade by them: transparency conditions written into public procurement of algorithmic systems, population-level impact assessment rather than individualised remedies \parencite{erdogan2024}, and statutory refusal of national-security carve-outs for municipal deployments. These are modest, actionable instruments — and the corpus suggests they address the box where it actually lives.

Third, \textit{treat `smart city' procurement as security policy}. The practitioner literature already does (§\ref{sec:theme1}): NATO-adjacent papers design the military use of civilian urban IoT; defence assessments treat adversary smart-city exports as strategic threats \parencite{atha2020,weber2023}. Municipalities buying integrated urban platforms are, on the evidence of this corpus, procuring battlespace infrastructure — whether or not they intend to. Democratic oversight that begins at the police procurement stage arrives years too late; it must begin at the platform contract.

Fourth, \textit{the warrant problem requires attention to the targeting side of dual-use}. Section~\ref{sec:dual-use-warrant} implies that humanitarian-law advocacy cannot cede the dual-use category to military lawyers. The emerging counter-forensic practice — using the same sensing stacks to document destruction that targeting uses to produce it \parencite{holail2024} — is the corpus's clearest existing model of contestation at this level, and it deserves policy support: open damage-assessment data, funded documentation mandates, and evidentiary standards for algorithmically generated targets.

\section{What the Map Cannot Say}
\label{sec:map-cannot-say}

The limitations registered in §\ref{sec:limitations} constrain these implications in ways worth restating substantively rather than procedurally. The review maps public discourse: the classified systems that would most directly test the thesis are absent by definition, and the five-record intersection may partly reflect that absence rather than a true gap in the world. The corpus skews Anglophone and theoretical; quantitative evaluation of algorithmic systems' battlefield performance is nearly absent (7\% of records). The Maybes dropped from the snowball may conceal relevant bridges. And the Palestine concentration of the destruction literature — paradigmatically productive but geographically narrow — means claims about algorithmic urbicide in general are built disproportionately on one case, however well-documented.

The conclusions in Chapter~6 are calibrated to these constraints: the dissertation claims to have mapped and explained a discursive junction, not to have demonstrated the operational reality of any specific targeting system.
